\chapter{Technische Grundlagen}
\abschnittsautor{M. Musterfrau}

Hier steht alles, was zum Verständnis der praktischen Umsetzung nötig ist --
und nur das. Jede Technologie, die ihr später verwendet, wird hier einmal
erklärt und belegt.

\section{Teilbereich A}

Beim ersten Auftreten schreibt ihr eine Abkürzung aus: \ac{API}. Danach
genügt die Kurzform: \ac{API}. Alle Abkürzungen sammelt ihr in
\texttt{diplomarbeit.tex} im Abkürzungsverzeichnis.

\subsection{Ein Unterabschnitt}

Zitiert wird im IEEE-Stil mit einer Nummer in eckigen
Klammern~\cite{beispiel-artikel}. Mehrere Quellen auf einmal gehen
auch~\cite{beispiel-buch,beispiel-online}.

\subsubsection{Eine Ebene tiefer}

Die vierte Ebene könnt ihr verwenden, aber sparsam. Sie wird nummeriert,
steht aber nicht mehr im Inhaltsverzeichnis. Tiefer gliedert ihr nicht.

\section{Teilbereich B}

Erklärt Technologien in eurer eigenen Formulierung. Wörtlich übernommene
Passagen gehören in Anführungszeichen und mit Seitenangabe belegt, sonst gilt
das als Plagiat.
